\section{结论与未来工作}

% ─────────────────────────────────────────────────────────
\subsection{结论}

本文提出了 \textbf{Camera-Music ControlNet}——一种将视频镜头语言方法论系统性迁移至音乐生成领域的创新框架。
面对 AI 音乐生成领域长期存在的“提示词工程代差”——视频生成已具备多维度镜头语言控制体系，
而音乐生成仍停留在粗粒度标签层面——我们识别出根本原因在于该领域缺乏一套 AI 可解析的、
人类可读的音乐 Prompt DSL，并以此为核心动机展开了系统性的技术攻关。

本文的主要贡献可概括为以下四点：
(\textit{i}) \textbf{六维 Token DSL 的形式化框架}——建立了视频镜头语言与音乐学术语之间的
九组结构性对应关系，并定义了覆盖 Harmony、Rhythm、Texture、Dynamics、Timbre 和 Articulation
六个正交维度的可扩展 Token 词汇表。零样本验证实验确认 MusicGen 预训练模型对五维度 Token
具有超出预期的响应度（音色维度频谱质心差 13.6$\times$，奏法维度 onset 10 vs 0）。
(\textit{ii}) \textbf{0.67\% 参数的 ControlNet-Adapter}——基于 Concatenation 注入法，
在完全冻结 MusicGen 解码器的前提下实现了六维精准控制。消融实验表明，
该适配器对和弦维度（Chord Accuracy 降幅 34.6\%）和节奏维度（Rhythm Obedience 降幅 53.7\%）
的依赖性最强，且六维度之间呈现显著的互补性。
(\textit{iii}) \textbf{热插拔混合专家架构}——基于显式 \texttt{[STYLE]} Token 路由机制
构建了三专家 MoE 系统，通过 Router 权重矩阵的行追加操作（Explicit Clone）实现了新风格的无损扩展，
古风专家额外引入了五声音阶偏置、民乐音色调制和装饰音特征三项结构先验。
(\textit{iv}) \textbf{零成本国风自动化数据工厂}——设计了 MIDI + 民乐 SF2 渲染 + 二胡单音化
+ 大语言模型自动标注的全自动数据管线，已产出 25 首二胡纯音渲染国风曲目，
为“人类可读的音乐 Prompt DSL”从概念到数据再到模型的完整技术链路提供了验证。

% ─────────────────────────────────────────────────────────
\subsection{未来工作}

本文的工作开启了若干值得深入探索的方向。
\textbf{在数据层面}，我们计划利用学校 A100/H100 算力集群，将国风种子集从当前的 25 首
扩展至 500--1000 首规模，并引入古筝、琵琶、笛子等更多民乐 SF2 音源，
形成覆盖中国民族乐器全谱系的自动化语料生产线。
同时，我们将通过腾讯 SongPrep-7B\cite{songprep} 对合成数据进行自动结构标注和质量校验，
以提升训练数据的标注精度。

\textbf{在架构层面}，热插拔 MoE 的可扩展性已得到数学证明和代码验证，
下一步计划将风格专家从当前的 3 个（古风/民谣/流行）扩展至 5--8 个，
覆盖摇滚、电子、爵士、R\&B 等更多垂直音乐风格，以实证框架的规模无关特性。
此外，我们将在云 GPU 集群上完成 ControlNet-Adapter 的完整训练和六维消融实验的全指标复现，
用真实训练曲线替代当前的模拟数据。

\textbf{在应用层面}，我们正在开发基于大语言模型的“音乐指挥家”Agent 前端——
用户仅需输入自然语言意图（如“写一首凄美的古风离别曲”），
LLM 自动完成歌词创作和 Token 编排，下游的 MusicGen+MoE 直接合成带控制的完整音频。
该工作流旨在让非音乐专业的普通用户也能通过自然语言实现专业级的音乐创作控制。